% C-23 — Reactive Inspection Task Limitations paragraph, redundancy pass
% Origin: intake/pending/redundancy_closing_audit.md (R2-01-derivative redundancy audit)
% File: sections/results.tex

% --- BEFORE (ssec:Limitations) ---
% \revblue{Regarding the Reactive Inspection Task: the platform performs no deliberate path
% planning, builds no map, and does not compute a global re-route when blocked --- obstacle
% traversal emerges from local BG/WTA dynamics rather than global trajectory planning. Coverage
% rate as a formal area-inspection metric was not evaluated, as it presupposes an explicit
% coverage planner outside this work's scope; task-completion indicators (success rate, final
% approach precision, time to first visual acquisition) are reported instead. Visual target
% pursuit is conditioned on the absence of a frontal obstacle, suspending during evasion and
% resuming automatically once cleared. Finally, in the obstacle-field scenario used for this
% experiment the target is not visible from the initial pose, so the reported acquisition
% metric characterizes first acquisition rather than re-acquisition after occlusion.}

% --- AFTER (ssec:Limitations) ---
\revblue{The Reactive Inspection Task's scope constraints ---absence of deliberate path planning
or mapping, the resulting inapplicability of a formal coverage-rate metric, and the first- rather
than re-acquisition semantics of $T_{acquisition}$--- are established alongside that experiment's
own results in Section~\ref{sssec:reactive_inspection} above and are not restated here.}

% Also added \label{sssec:reactive_inspection} right after
% \subsubsection{Reactive Inspection Task} (previously unlabeled) so the pointer above resolves.

% Rationale: all four facts in the "before" paragraph (no path planning/map/re-route; coverage
% rate not reported + why; visual pursuit gated on frontal-obstacle absence; first- vs.
% re-acquisition) were already stated, with the same reasoning, in the Reactive Inspection Task's
% own results paragraphs (results.tex, "To directly address..." / "The target was first visually
% acquired..." paragraphs). The Limitations copy added zero new information -- pure restatement.
% Condensed to an establishing-mention + cross-reference, matching R2-01's pattern for other
% cross-section repeats.
